以下是我和chatgpt讨论的。不过它目前不知道我们的nn架构，所以我先写一个比较抽象的版本，后续可以根据实际架构细化。

这里面定义清楚了lossfunction和gradient path。后续我们可以在这个基础上添加一些细节，比如memory的输入输出，或者是一些特殊的loss项。

另外，里面也写了一些nonlinear system的例子。我们可以先从一个简单的scalar nonlinear system开始测试，看看能不能学会稳定控制。之后再逐步增加复杂度，比如多维系统，或者是有一些不确定性的系统。

\documentclass[11pt]{ctexart}

\usepackage{amsmath, amssymb, amsthm, mathtools}
\usepackage{enumitem}
\usepackage{geometry}
\usepackage{hyperref}
\geometry{margin=1in}

\newtheorem{problem}{Problem}
\newtheorem{experiment}{Experiment}
\newtheorem{remark}{Remark}

\title{Experiment Plan: Online Neural Control Without Explicit Equilibrium Computation}
\author{}
\date{}

\begin{document}

\maketitle

\section{Motivation}

We want to test whether a neural network controller can regulate an unknown nonlinear dynamical system without explicitly computing its steady state, without linearizing the dynamics, and without designing a complicated reward function.

The guiding idea is not classical stabilization of a prescribed equilibrium \(x^*\). Instead, the controller interacts with the environment, observes the consequence of its action, and updates itself using a simple local loss. The controller is expected to discover a relatively stable, low-energy regime by exploration and online learning.

The central question is:
\[
\text{Can a neural controller learn stabilizing behavior from local feedback alone?}
\]

\section{Basic Dynamical System}

We begin with a discrete-time controlled system
\[
x_t = F(x_{t-1},u_t),
\]
where \(x_t\in \mathbb R^n\) is the state after the input \(u_t\in \mathbb R^m\) has acted on the system.

This convention is important: \(x_t\) is the result generated by \(u_t\). Therefore, the instantaneous loss
\[
\ell_t = \|x_t\|^2 + \rho \|u_t\|^2
\]
does depend on \(u_t\) through the map \(F\).

For a continuous-time system
\[
\dot x = f(x,u),
\]
we use the Euler discretization
\[
x_t = x_{t-1} + \Delta t\, f(x_{t-1},u_t).
\]

\section{Minimal Toy Environment}

The first toy example is a scalar nonlinear unstable system:
\[
\dot x = x + x^3 + u.
\]
The discrete-time version is
\[
x_t = x_{t-1} + \Delta t\left(x_{t-1}+x_{t-1}^3+u_t\right).
\]

This system is useful because without control the term \(x+x^3\) tends to push the state away from the origin. A successful controller should learn to counteract this growth.

\section{Neural Controller}

The controller is a neural network
\[
u_t = \pi_\theta(o_{t-1}),
\]
where \(o_{t-1}\) is the observation available before applying \(u_t\).

In the simplest experiment,
\[
o_{t-1}=x_{t-1}.
\]
Thus,
\[
u_t = \pi_\theta(x_{t-1}).
\]

To avoid extremely large inputs, we may use a bounded output layer:
\[
u_t = u_{\max}\tanh(\mathrm{MLP}_\theta(x_{t-1})).
\]

\section{Local Instantaneous Loss}

We do not use a full trajectory loss from \(t=1\) to \(T\). Instead, the controller is updated locally at each time step using
\[
\ell_t
=
\|x_t\|^2+\rho\|u_t\|^2.
\]

Here \(\rho>0\) balances state regulation and control effort.

The interpretation is simple:
\[
\|x_t\|^2
\quad \text{penalizes large state magnitude,}
\]
and
\[
\rho\|u_t\|^2
\quad \text{penalizes excessive control effort.}
\]

We use addition rather than multiplication. The product loss
\[
\|x_t\|\,\|u_t\|
\]
is not appropriate as the main loss because it becomes zero whenever either \(x_t=0\) or \(u_t=0\), and therefore does not separately penalize large states and large controls.

\section{Gradient Propagation}

Since \(x_t=F(x_{t-1},u_t)\), the local loss depends on \(u_t\). If \(F\) is differentiable, then
\[
\frac{\partial \ell_t}{\partial u_t}
=
2\rho u_t
+
\left(\frac{\partial x_t}{\partial u_t}\right)^\top 2x_t.
\]

Then the neural controller is updated by
\[
\nabla_\theta \ell_t
=
\frac{\partial \ell_t}{\partial u_t}
\frac{\partial u_t}{\partial \theta}.
\]

The gradient path is
\[
\ell_t
\longrightarrow
x_t
\longrightarrow
u_t
\longrightarrow
\theta.
\]

In implementation, for one-step local learning, we detach the previous state \(x_{t-1}\) from the computational graph. Thus the update only accounts for the effect of the current action \(u_t\) on the current resulting state \(x_t\), rather than backpropagating through the whole past trajectory.

\section{Experiment 1: Differentiable Environment}

\begin{experiment}[Known differentiable environment]
Use the true differentiable simulator
\[
x_t = F(x_{t-1},u_t)
\]
inside PyTorch or JAX.

At each time step:
\[
u_t=\pi_\theta(x_{t-1}),
\]
\[
x_t=F(x_{t-1},u_t),
\]
\[
\ell_t=\|x_t\|^2+\rho\|u_t\|^2.
\]

Then update
\[
\theta \leftarrow \theta-\eta\nabla_\theta \ell_t.
\]

This experiment verifies that the basic gradient pathway works when the environment is differentiable.
\end{experiment}

\subsection*{Algorithm}

\[
\begin{aligned}
&\text{Given initial state }x_0, \text{ learning rate }\eta, \text{ time step }\Delta t.\\
&\text{For }t=1,2,\dots:\\
&\quad \text{detach }x_{t-1}.\\
&\quad u_t=\pi_\theta(x_{t-1}).\\
&\quad x_t=x_{t-1}+\Delta t\left(x_{t-1}+x_{t-1}^3+u_t\right).\\
&\quad \ell_t=x_t^2+\rho u_t^2.\\
&\quad \theta\leftarrow \theta-\eta\nabla_\theta \ell_t.
\end{aligned}
\]

\section{Experiment 2: Unknown Environment with Learned Predictor}

In the second experiment, the controller is not allowed to backpropagate through the true environment. Instead, we introduce a learned predictor
\[
\hat x_t=P_\phi(x_{t-1},u_t).
\]

The real environment still produces
\[
x_t=F(x_{t-1},u_t),
\]
but the controller receives gradients through the differentiable predictor.

The prediction error is
\[
e_t=x_t-\hat x_t.
\]

The predictor is trained by
\[
\ell_t^{\rm pred}
=
\|x_t-\hat x_t\|^2.
\]

The controller is trained using
\[
\ell_t^{\rm ctrl}
=
\|\hat x_t\|^2+\rho\|u_t\|^2.
\]

The controller gradient path is now
\[
\ell_t^{\rm ctrl}
\longrightarrow
\hat x_t
\longrightarrow
u_t
\longrightarrow
\theta.
\]

The predictor receives the model-learning update
\[
\phi \leftarrow \phi-\eta_p\nabla_\phi \ell_t^{\rm pred},
\]
and the controller receives the control update
\[
\theta \leftarrow \theta-\eta_c\nabla_\theta \ell_t^{\rm ctrl}.
\]

\section{Experiment 3: Add Memory and Prediction Error Feedback}

After the one-step controller works, we introduce memory:
\[
m_t=G_\theta(m_{t-1},x_{t-1},u_t,e_t).
\]

The controller becomes
\[
u_t=\pi_\theta(m_{t-1},x_{t-1},e_{t-1}).
\]

The predictor becomes
\[
\hat x_t=P_\phi(m_{t-1},x_{t-1},u_t).
\]

The closed-loop system is
\[
\begin{aligned}
u_t &= \pi_\theta(m_{t-1},x_{t-1},e_{t-1}),\\
\hat x_t &= P_\phi(m_{t-1},x_{t-1},u_t),\\
x_t &= F(x_{t-1},u_t),\\
e_t &= x_t-\hat x_t,\\
m_t &= G_\theta(m_{t-1},x_{t-1},u_t,e_t).
\end{aligned}
\]

This version tests whether prediction error can serve as an internal feedback signal for adaptive control.

\section{Experiment 4: Slowly Changing Environment}

To test adaptation, we allow the environment parameter to vary:
\[
\dot x = a x+x^3+u,
\]
where \(a\) changes across episodes or slowly changes over time.

For example,
\[
a\in [0.5,2.0].
\]

The discrete dynamics are
\[
x_t=x_{t-1}+\Delta t\left(a x_{t-1}+x_{t-1}^3+u_t\right).
\]

The controller does not observe \(a\). It only observes \(x_{t-1}\), \(u_t\), and possibly the prediction error \(e_t\). The purpose is to test whether memory and prediction error help the controller adapt to different environments.

\section{Evaluation Metrics}

For each experiment, record and plot:

\[
x_t \quad \text{over time},
\]
\[
u_t \quad \text{over time},
\]
\[
\ell_t=\|x_t\|^2+\rho\|u_t\|^2 \quad \text{over time},
\]
and, when using a predictor,
\[
e_t=x_t-\hat x_t \quad \text{over time}.
\]

Successful behavior means:
\[
x_t \text{ remains bounded and tends to a small region,}
\]
\[
u_t \text{ remains bounded and not excessively large,}
\]
\[
\ell_t \text{ decreases or remains small,}
\]
and, for predictor-based control,
\[
e_t \to 0 \quad \text{or remains small.}
\]

\section{Expected Outcomes}

The expected result of Experiment 1 is that the neural controller learns a feedback-like policy that suppresses the unstable nonlinear growth of
\[
\dot x=x+x^3+u.
\]

The expected result of Experiment 2 is that the controller can still learn stabilizing behavior when gradients are propagated through a learned predictor rather than through the true environment.

The expected result of Experiment 3 is that memory and prediction error improve adaptation and robustness.

The expected result of Experiment 4 is that the memory-based controller can adapt to changing dynamics without explicitly identifying the parameter \(a\).

\section{Important Conceptual Point}

This experiment does not claim classical stabilization to a prescribed equilibrium \(x^*\). Instead, the goal is to study a weaker and more biologically inspired notion:

\[
\textit{Can a neural controller discover a stable, low-energy regime through local online learning?}
\]

In this sense, the controller does not compute a steady state, does not linearize the system, and does not use a handcrafted complicated reward. It only receives local feedback through the simple instantaneous loss
\[
\ell_t=\|x_t\|^2+\rho\|u_t\|^2.
\]

\end{document}